\section{A faster raw oracle for sub-Gaussian labels}
\label{sec:raw-truncation-transfer}
A native uniform-output tail also permits a second statistical argument.
It uses a bounded prior only as an auxiliary object in the proof. The
actual prior, Gaussian likelihood, and posterior mean remain unbounded.

\begin{lemma}[Bounded predictions with arbitrary finite-variance labels]
\label{lem:bounded-finitevariance-loo}
Let $\pi$ be any prior on functions taking values in $[-M,M]$, with
$M\ge1$, and use temperature $T\ge256M^2$. For IID observations with
$\E Y^2<\infty$, the ordinary posterior mean trained on all $n$
observations obeys
\[
 \E R(\bar f_n)\le\int R(f)\,d\rho
       +\frac{T\KL(\rho\Vert\pi)}{n+1}
       +4\E[Y^2\mathbf1_{\{|Y|>M\}}]
\]
for every deterministic comparison law $\rho$.
\end{lemma}
\begin{proof}
For a distribution $q$ on $[-M,M]$ with mean $m$ and variance $v$,
let $p_y(du)\propto e^{-(u-y)^2/T}q(du)$. If $|y|\le M$, the
covariance identity and the mean-value theorem give
\[
 \E_{p_y}(u-y)^2\ge(m-y)^2+
 \left[1-\frac{4(M+|y|)^2}{T}
           e^{4M(M+|y|)/T}\right]v\ge(m-y)^2.
\]
For completeness, write $L=(u-y)^2$, $w=e^{-L/T}$. Then
$\E_{p_y}L=\E_qL+\operatorname{Cov}_q(L,w)/\E_qw$.
The two-copy covariance identity, $|L(u)-L(u')|\le
2(M+|y|)|u-u'|$, and
$|w(u)-w(u')|\le e^{-\min L/T}|L(u)-L(u')|/T$
bound the possible decrease by
$4(M+|y|)^2e^{\operatorname{range}(L)/T}v/T$.
The displayed conservative range bound proves the inequality.
If $|y|>M$, the squared-loss range over $[-M,M]$ is exactly
$4M|y|\le4y^2$, so the same inequality holds with additive error
$4y^2\mathbf1_{\{|y|>M\}}$.

Apply this scalar assertion to each prediction from the posterior
trained on the other $n$ points of an IID sample of size $n+1$.
Exchangeability identifies its expected mean loss with the left side.
After the extra likelihood factor, each term uses the full posterior
on $n+1$ observations. Its expected empirical risk is at most
$\int R\,d\rho+T\KL(\rho\Vert\pi)/(n+1)$ by the variational
inequality. The prior, $M$, and $T$ are held fixed at their
$n$-specific values throughout this leave-one-out argument.
\end{proof}

\begin{theorem}[Sub-Weibull priors and a raw sub-Gaussian-label oracle]
\label{thm:raw-tail-oracle}
Suppose a centered function prior satisfies, uniformly over a family
of finite architectures,
\[
 \|U(f)\|_{L^p(\mu)}\le A p^r\quad(p\ge1),\qquad r>0,
 \qquad U(f)=\sup_x|f(x)|.
\]
Assume $\E e^{Y^2/K_Y^2}\le2$. There are constants $c_0,c_1,c_2$,
depending only on $A,r,K_Y$, such that the following holds. Put
\[
 b=\frac1{2r+1},\qquad q_n=c_0n^b,\qquad
 M_n=eAq_n^r,\qquad T_n=256M_n^2=c_1n^{1-b}.
\]
Use the actual unrestricted prior and the raw Gaussian-likelihood
posterior with temperature $T_n$. If a deterministic comparison
function $g$ obeys $\|g\|_\infty\le M_n/2$, and a law $\rho$
satisfies
\[
 \E_\rho\|f-g\|_\infty^2
       \le\frac{(n+1)^{-4}}{1+\|g\|_\infty^2},
 \qquad\KL(\rho\Vert\mu)\le K_n(g),
\]
then its ordinary raw posterior mean satisfies
\begin{equation}
 \E R(\bar f_{S_n})\le R(g)+c_2 n^{-b}K_n(g)
                                  +c_2(n+1)^{-2}.
 \label{eq:raw-tail-oracle}
\end{equation}
The prior is not truncated, and no output link is used in the learner.
\end{theorem}
\begin{proof}
Let $E_n=\{U\le M_n\}$. Markov's inequality at moment $q_n$ gives
$\mu(E_n^c)\le e^{-q_n}$. The auxiliary prior
$\mu_n=\mu(\cdot\mid E_n)$ has outputs bounded by $M_n$.
The comparison's probability of $E_n^c$ is at most
$4(n+1)^{-4}/M_n^2$, so it is at most $1/2$ after increasing $c_0$.
Condition the comparison on $E_n$. The KL chain rule against the
partition $(E_n,E_n^c)$ gives
\[
 \KL(\rho(\cdot\mid E_n)\Vert\mu_n)
        \le\KL(\rho\Vert\mu)/\rho(E_n)\le2K_n(g).
\]
Its mean squared approximation error increases by at most two.
Cauchy--Schwarz consequently bounds its risk by
$R(g)+C_{K_Y}(n+1)^{-2}$. Lemma~\ref{lem:bounded-finitevariance-loo}
now bounds the auxiliary posterior-mean risk by
\[
 R(g)+\frac{2T_nK_n(g)}{n+1}
   +C_{K_Y}(n+1)^{-2}
   +4\E[Y^2\mathbf1_{\{|Y|>M_n\}}].
\]
The last term is at most $C_{A,r,K_Y}(n+1)^{-2}$ after increasing
$c_0$, by the sub-Gaussian label tail.

It remains to compare the two actual prediction rules. Let $m_S$ be
the unrestricted posterior mean and $m_{S,E}$ the mean under
$\mu_n$. The latter posterior is exactly the former conditioned on
$E_n$. Centering of the native prior and the bound
$V:=\sup_x\E_\mu f(x)^2\le A^22^{2r}$ imply
\[
 Z_S\ge\exp[-(nV+\sum_iY_i^2)/T_n].
\]
Therefore, uniformly over inputs,
\begin{align*}
 |m_S-m_{S,E}|
 &\le Z_S^{-1}\E_\mu[(U+M_n)\mathbf1_{E_n^c}]\\
 &\le C_A(1+M_n)
     \exp\{(nV+\sum_iY_i^2)/T_n-q_n/2\}=:D_S.
\end{align*}
The second inequality uses Cauchy--Schwarz and the fixed native second
moment. For $T_n\ge2K_Y^2$, Lyapunov's inequality gives
\[
 \E e^{2\sum_iY_i^2/T_n}
 \le\exp(2nK_Y^2\log2/T_n).
\]
Choose $c_0$ large enough that
$q_n\ge8n(V+K_Y^2)/T_n+C\log(n+1)+C\log(1+M_n)$ for every $n$.
This is possible: both $q_n$ and $n/T_n$ are constant multiples of
$n^b$, and increasing $c_0$ increases the former while decreasing
the latter. It follows that
$\E[2(M_n+\E|Y|)D_S+D_S^2]\le C(n+1)^{-2}$.
Finally
\[
 R(m_S)\le R(m_{S,E})+
                   2(M_n+\E|Y|)D_S+D_S^2.
\]
Combining the estimates proves \eqref{eq:raw-tail-oracle}. The
auxiliary truncation has disappeared from the conclusion and from the
specified learning algorithm.
\end{proof}

\begin{corollary}[A comparison already inside the analytical cutoff]
\label{cor:raw-supported-comparison}
Under the prior and label assumptions of
Theorem~\ref{thm:raw-tail-oracle}, let $T_n,M_n$ be its deterministic
temperature and cutoff. If a sample-independent law $\rho$ is supported
on $\{U\le M_n\}$, then
\[
 \E R(\bar f_{S_n})\le\int R(f)\,d\rho(f)
       +\frac{T_n\KL(\rho\Vert\mu)}{n+1}+C(n+1)^{-2}.
\]
No uniform approximation to a particular comparator is needed; a
population-risk comparison is enough.
\end{corollary}
\begin{proof}
Use $\rho$ directly in the auxiliary bounded-prior theorem. Since it
is already supported on $E_n$,
$\KL(\rho\Vert\mu(\cdot\mid E_n))=
 \KL(\rho\Vert\mu)+\log\mu(E_n)\le\KL(\rho\Vert\mu)$.
The label tail and actual-versus-auxiliary posterior transfer are
exactly those in the preceding proof, independent of the comparison.
\end{proof}
